% Part: history
% Chapter: biographies
% Section: bertrand-russell

\documentclass[../../../include/open-logic-section]{subfiles}

\begin{document}

\olfileid{his}{bio}{rus} 

\olsection{બર્ટ્રાન્ડ રસેલ}

\olphoto{russell-bertrand}{બર્ટ્રાન્ડ રસેલ}

બર્ટ્રાન્ડ રસેલને આધુનિક
વિશ્લેષણાત્મક તત્વજ્ઞાનના
સ્થાપકોમાંના એક ગણવામાં
આવે છે. 18 મે 1872ના રોજ
જન્મેલા રસેલ તત્વજ્ઞાન અને
તર્કશાસ્ત્રના કાર્ય માટે
જ જાણીતા ન હતા; તેમણે
વિવિધ વિષયો પર લોકપ્રિય
પુસ્તકો પણ લખ્યાં હતાં.
જીવનભર તેઓ ઉત્સાહી
રાજકીય કાર્યકર પણ રહ્યા.

રસેલનો જન્મ વેલ્સના
મૉનમથશાયરના ટ્રેલેકમાં
થયો હતો. તેમનાં માતાપિતા
બ્રિટિશ અમીરવર્ગનાં હતાં.
તેઓ મુક્તચિંતક હતાં અને
તે સમયના બોસ્ટનના
ક્રાંતિકારી વિચારવાળાઓ
સાથે પણ મિત્રતા બાંધી
હતી. દુર્ભાગ્યે રસેલ
નાના હતા ત્યારે જ
તેમનાં માતાપિતાનું
અવસાન થયું અને
તેમને દાદા-દાદી
સાથે રહેવા મોકલાયા.
ત્યાં તેમનો ઉછેર
ધાર્મિક વાતાવરણમાં
થયો, જ્યારે તેમના
માતાપિતા કોઈ પણ
ભોગે એવું ટાળવા
માગતા હતા. તેમની
દાદી નૈતિકતાની
બધી બાબતોમાં બહુ
કડક હતાં. કિશોરાવસ્થામાં
મોટાભાગે ખાનગી
શિક્ષકો તેમને ઘરે
જ ભણાવતા.

વિશ્લેષણાત્મક તત્વજ્ઞાનમાં,
ખાસ કરીને તર્કશાસ્ત્રમાં,
રસેલનો પ્રભાવ બહુ મોટો
છે. તેમણે કેમ્બ્રિજની
ટ્રિનિટી કૉલેજમાં
ગણિત અને તત્વજ્ઞાન
ભણ્યાં; ત્યાં ગણિતશાસ્ત્રી
અને ફિલસૂફ આલ્ફ્રેડ
નૉર્થ વ્હાઇટહેડનો
તેમના પર પ્રભાવ
પડ્યો. 1910માં રસેલ
અને વ્હાઇટહેડે
\emph{Principia Mathematica}નો
પ્રથમ ખંડ પ્રકાશિત
કર્યો, જેમાં તેમણે
ગણિતને તર્કશાસ્ત્રમાં
ઘટાડી શકાય છે એવો
દૃષ્ટિકોણ રજૂ કર્યો.
પછી તેમણે સેંકડો
પુસ્તકો, નિબંધો અને
રાજકીય પત્રિકાઓ
પ્રકાશિત કરી. 1950માં
તેમને સાહિત્યનું
નોબેલ પારિતોષિક મળ્યું.

રસેલ રાજકારણ અને
સામાજિક ચળવળોમાં
ઊંડે જોડાયેલા હતા.
પ્રથમ વિશ્વયુદ્ધ
દરમિયાન શાંતિવાદી
પ્રવૃત્તિઓ અને
વિરોધપ્રદર્શનને
કારણે તેમની
ધરપકડ થઈ અને
તેમને છ મહિના
જેલમાં રાખવામાં
આવ્યા. જેલમાં
તેઓ વાંચી-લખી
શકતા હતા અને
પોતાને એ અનુભવ
“ખાસ્સો સુખદ”
લાગ્યો હોવાનું
કહેતા. તેઓ જીવનભર
શાંતિવાદી રહ્યા;
1961માં પરમાણુ
નિઃશસ્ત્રીકરણની
સભામાં હાજરી
આપવા બદલ તેમને
ફરી કેદ કરવામાં
આવ્યા. 1948માં
તેઓ વિમાન દુર્ઘટનામાંથી
પણ બચ્યા; તેમાં
માત્ર ધૂમ્રપાન
માટેના વિભાગમાં
બેઠેલા લોકો જ
બચ્યા હતા. તેથી
રસેલ કહેતા કે
પોતાનો જીવ
ધૂમ્રપાનને આભારી
છે. તેમણે ચાર
વાર લગ્ન કર્યાં,
પણ લગ્નબાહ્ય
સંબંધો માટે પણ
તેમની ખ્યાતિ હતી.
2 ફેબ્રુઆરી 1970ના
રોજ 97 વર્ષની
ઉંમરે વેલ્સના
પેનરિનડેઉડ્રાઇથમાં
તેમનું અવસાન થયું.

\begin{reading}
રસેલે 1872--1967નો
પોતાનો જીવનકાળ
આવરી લેતી ત્રણ
ભાગની આત્મકથા
લખી હતી
\citep{Russell1967,Russell1968,Russell1969}.
મૅકમાસ્ટર યુનિવર્સિટીમાં
બર્ટ્રાન્ડ રસેલ સંશોધન
કેન્દ્ર તેમની દસ્તાવેજી
સામગ્રીનું ભંડાર છે.
તેમનાં એકત્રિત લખાણોના
ખંડો, તેમાં શોધ કરી
શકાય તેવી સૂચિઓ,
અને સંગ્રહ-પ્રકલ્પો
વિશે જાણકારી માટે
કેન્દ્રની વેબસાઇટ
\citet{Duncan2015} જુઓ.
રસેલનો
\emph{On
  Denoting}
\citep{Russell1905}
વીસમી સદીના
વિશ્લેષણાત્મક
તત્વજ્ઞાનનો ઉત્તમ
લેખ છે.

રસેલ અંગેની
The Stanford Encyclopedia of Philosophyની
નોંધમાં તેમના
ઇચ્છા અને રાજકીય
સિદ્ધાંત વિશેના
વક્તવ્યની ધ્વનિ-ક્લિપ્સ
છે \citep{Irvine2015}.
તેમની ઘણી વિડિયો
મુલાકાતો ઑનલાઇન
ઉપલબ્ધ છે. દાખલા
તરીકે, ધૂમ્રપાન
અને વિમાન દુર્ઘટના
વિશે તેમને બોલતા
જોવા \citet{RussellND}
જુઓ. તેમની કેટલીક
કૃતિઓ, જેમાં
\emph{Introduction to Mathematical Philosophy}
પણ છે, \citet{LibriVoxND}
પર મફત શ્રાવ્ય
પુસ્તકો તરીકે
ઉપલબ્ધ છે.
\end{reading}

\end{document}
